\documentclass[11pt]{article}
\usepackage{amsmath, amssymb, geometry, graphicx, enumerate}
\geometry{a4paper, margin=1in}

\title{\textbf{Homework 1.2: Row Reduction and Echelon Forms}}
\author{Tashfeen Omran}
\date{September 2026}

\begin{document}

\maketitle
\tableofcontents
\vspace{1cm}
\hrule
\vspace{1cm}

\section{Part 1: Comprehensive Introduction, Context, and Prerequisites}

\subsection{Core Concepts}
In the previous section, we learned that elementary row operations can be used to solve systems of linear equations. In Section 1.2, we formalize this process into a systematic algorithm called \textbf{row reduction}. The goal of row reduction is to transform a complex matrix into a simplified, staircase-like structure that makes finding the solution trivial. We define two crucial forms:

\textbf{1. Echelon Form (or Row Echelon Form, REF):}
A matrix is in echelon form if it satisfies three conditions:
\begin{enumerate}
    \item All nonzero rows are above any rows of all zeros.
    \item Each leading entry (the leftmost nonzero entry) of a row is in a column to the right of the leading entry of the row above it.
    \item All entries in a column below a leading entry are zeros.
\end{enumerate}

\textbf{2. Reduced Echelon Form (or Reduced Row Echelon Form, RREF):}
A matrix is in reduced echelon form if it satisfies the three conditions above, plus two more:
\begin{enumerate}
    \setcounter{enumi}{3}
    \item The leading entry in each nonzero row is $1$.
    \item Each leading $1$ is the \textbf{only} nonzero entry in its column (i.e., there are zeros both below \textit{and} above the leading $1$).
\end{enumerate}

\textbf{Pivots and Variables:}
\begin{itemize}
    \item A \textbf{pivot position} in a matrix $A$ is a location that corresponds to a leading $1$ in the reduced echelon form of $A$.
    \item A \textbf{pivot column} is a column of $A$ that contains a pivot position.
    \item Variables corresponding to pivot columns are called \textbf{basic variables}.
    \item Variables corresponding to non-pivot columns are called \textbf{free variables}.
\end{itemize}

\subsection{Intuition and Derivation}
The algorithm to reach RREF is divided into two phases. The \textbf{forward phase} works from left to right, top to bottom, creating zeros below pivots to reach echelon form. This is mathematically equivalent to the "elimination" method in algebra. The \textbf{backward phase} works from right to left, bottom to top, creating zeros above the pivots and scaling them to $1$. This is the matrix equivalent of "back-substitution."

When a system has free variables, it means we do not have enough unique information to pin down a single point. The free variables act as parameters (often denoted $t, s, \dots$) that can take on any real value, tracing out a line, plane, or hyperplane of infinite solutions.

\subsection{Historical Context and Motivation}
While Carl Friedrich Gauss popularized the elimination method (now called Gaussian elimination) in the early 1800s for astronomical computations, it was Wilhelm Jordan in 1888 who extended the algorithm to what we now call \textbf{Gauss-Jordan elimination}. Jordan realized that by continuing the elimination process upwards (the backward phase) to achieve reduced echelon form, one could read the solutions directly from the matrix without doing algebraic back-substitution by hand. This algorithmic approach is the bedrock of modern computational linear algebra in computer science.

\subsection{Key Formulas and Theorems}
\textbf{Theorem 1: Uniqueness of the Reduced Echelon Form}
Each matrix is row equivalent to one and \textbf{only one} reduced echelon matrix. (Note: The standard echelon form is \textit{not} unique, as different sequences of row operations can yield different non-reduced echelon forms, but they will always eventually simplify to the exact same RREF).

\textbf{Theorem 2: Existence and Uniqueness Theorem}
A linear system is consistent if and only if the rightmost column of the augmented matrix is \textit{not} a pivot column. If a linear system is consistent, the solution set contains either:
\begin{enumerate}
    \item A unique solution (when there are no free variables).
    \item Infinitely many solutions (when there is at least one free variable).
\end{enumerate}

\subsection{Prerequisites}
To master this section, you must be comfortable with the three elementary row operations (replacement, interchange, scaling) from Section 1.1, and you must have a sharp eye for identifying leading entries and reading matrices left-to-right.

\newpage
\section{Part 2: Detailed Homework Solutions}

\subsection*{Problem 1.2.2}
\textbf{Question:} Determine which matrices are in reduced echelon form and which others are only in echelon form.
\[ a. \begin{bmatrix} 1 & 0 & 0 & 0 \\ 0 & 1 & 0 & 0 \\ 0 & 0 & 1 & 1 \end{bmatrix} \quad b. \begin{bmatrix} 0 & 1 & 1 & 1 & 1 \\ 0 & 0 & 1 & 1 & 1 \\ 0 & 0 & 0 & 0 & 1 \\ 0 & 0 & 0 & 0 & 0 \end{bmatrix} \quad c. \begin{bmatrix} 0 & 0 & 0 & 0 \\ 1 & 3 & 0 & 0 \\ 0 & 0 & 1 & 0 \\ 0 & 0 & 0 & 1 \end{bmatrix} \]

\textbf{Solution:}
\begin{itemize}
    \item \textbf{Matrix a:} The leading entries of the nonzero rows are all $1$. Each leading $1$ is further to the right than the one above it. The columns containing the leading $1$s (columns 1, 2, and 3) have zeros everywhere else in those columns. Therefore, it satisfies all five conditions. \\
    \textbf{Answer a: reduced echelon form.}
    
    \item \textbf{Matrix b:} The leading entries are $1$s, they move to the right as we go down, and all rows of zeros are at the bottom. The entries below the leading $1$s are all zero. This means it is in echelon form. However, look at the leading $1$ in row 2 (column 3). The entry \textit{above} it is a $1$, not a zero. Because leading entries are not the only nonzero entries in their columns, it is not reduced.\\
    \textbf{Answer b: echelon form only.}
    
    \item \textbf{Matrix c:} The very first row is a row of all zeros. However, there are nonzero rows beneath it. This violates the very first rule of echelon forms: all nonzero rows must be above any rows of all zeros.\\
    \textbf{Answer c: neither.}
\end{itemize}


\subsection*{Problem 1.2.4}
\textbf{Question:} Row reduce the matrix to reduced echelon form. Identify the pivot positions in the original matrix, and list the pivot columns.
\[ \begin{bmatrix} 1 & 2 & 4 & -9 \\ 2 & 4 & 5 & -12 \\ 4 & 5 & 4 & -15 \end{bmatrix} \]

\textbf{Solution:}
Let's row reduce the matrix to find the pivots.
\[ \underrightarrow{R_2 \leftarrow R_2 - 2R_1, \ R_3 \leftarrow R_3 - 4R_1} \begin{bmatrix} 1 & 2 & 4 & -9 \\ 0 & 0 & -3 & 6 \\ 0 & -3 & -12 & 21 \end{bmatrix} \]
To maintain the staircase pattern, we must swap $R_2$ and $R_3$.
\[ \underrightarrow{R_2 \leftrightarrow R_3} \begin{bmatrix} 1 & 2 & 4 & -9 \\ 0 & -3 & -12 & 21 \\ 0 & 0 & -3 & 6 \end{bmatrix} \]
The matrix is now in echelon form. The leading entries are in columns 1, 2, and 3. Therefore, \textbf{Columns 1, 2, and 3 are the pivot columns.}
A pivot position is a position $(i, j)$ in the matrix that corresponds to a leading entry in an echelon form. Since our pivots are in columns 1, 2, and 3 for rows 1, 2, and 3 respectively, the pivot positions are $(1,1)$, $(2,2)$, and $(3,3)$. 
In the original matrix, these positions contain the numbers $1, 4,$ and $4$.
\\
\textbf{Part 2 Answer:} Option B is correct because it correctly highlights the pivot positions corresponding to the leading entries. \\
\textbf{Part 3 Answer:} Pivot columns are Column 1, Column 2, and Column 3.


\subsection*{Problem 1.2.11}
\textbf{Question:} Find the general solution of the system whose augmented matrix is shown.
\[ \begin{bmatrix} 3 & -2 & 4 & 0 \\ -6 & 4 & -8 & 0 \\ -12 & 8 & -16 & 0 \end{bmatrix} \]

\textbf{Solution:}
Row reduce the matrix. Notice that row 2 and row 3 are multiples of row 1.
\[ \underrightarrow{R_2 \leftarrow R_2 + 2R_1, \ R_3 \leftarrow R_3 + 4R_1} \begin{bmatrix} 3 & -2 & 4 & 0 \\ 0 & 0 & 0 & 0 \\ 0 & 0 & 0 & 0 \end{bmatrix} \]
Divide $R_1$ by $3$ to achieve reduced echelon form:
\[ \begin{bmatrix} 1 & -2/3 & 4/3 & 0 \\ 0 & 0 & 0 & 0 \\ 0 & 0 & 0 & 0 \end{bmatrix} \]
The only pivot is in column 1. Therefore, $x_1$ is a basic variable, while $x_2$ and $x_3$ are free variables. 
Write the equation for the first row:
\[ 1x_1 - \frac{2}{3}x_2 + \frac{4}{3}x_3 = 0 \]
Solve for the basic variable $x_1$:
\[ x_1 = \frac{2}{3}x_2 - \frac{4}{3}x_3 \]
\\
\textbf{Answer:} C. $x_1 = 2/3x_2 - 4/3x_3$, $x_2$ is free, $x_3$ is free.


\subsection*{Problem 1.2.13}
\textbf{Question:} Find the general solution of the system whose augmented matrix is given below.
\[ \begin{bmatrix} 1 & -3 & 0 & -1 & 0 & -5 \\ 0 & 1 & 0 & 0 & -2 & 7 \\ 0 & 0 & 0 & 1 & 7 & 2 \\ 0 & 0 & 0 & 0 & 0 & 0 \end{bmatrix} \]

\textbf{Solution:}
The matrix is in echelon form, but not \textit{reduced} echelon form because there are non-zero entries above the pivots. The pivots are in columns 1, 2, and 4.
First, we use the pivot in $R_3$ (which is $1$ in column 4) to eliminate the $-1$ above it in $R_1$.
\[ \underrightarrow{R_1 \leftarrow R_1 + R_3} \begin{bmatrix} 1 & -3 & 0 & 0 & 7 & -3 \\ 0 & 1 & 0 & 0 & -2 & 7 \\ 0 & 0 & 0 & 1 & 7 & 2 \\ 0 & 0 & 0 & 0 & 0 & 0 \end{bmatrix} \]
Next, we use the pivot in $R_2$ (which is $1$ in column 2) to eliminate the $-3$ above it in $R_1$.
\[ \underrightarrow{R_1 \leftarrow R_1 + 3R_2} \begin{bmatrix} 1 & 0 & 0 & 0 & (7 + 3(-2)) & (-3 + 3(7)) \\ 0 & 1 & 0 & 0 & -2 & 7 \\ 0 & 0 & 0 & 1 & 7 & 2 \\ 0 & 0 & 0 & 0 & 0 & 0 \end{bmatrix} \]
Simplify the calculations in $R_1$: $7 - 6 = 1$, and $-3 + 21 = 18$.
\[ \begin{bmatrix} 1 & 0 & 0 & 0 & 1 & 18 \\ 0 & 1 & 0 & 0 & -2 & 7 \\ 0 & 0 & 0 & 1 & 7 & 2 \\ 0 & 0 & 0 & 0 & 0 & 0 \end{bmatrix} \]
The pivot columns are 1, 2, and 4. So $x_1, x_2, x_4$ are basic. Columns 3 and 5 lack pivots, making $x_3$ and $x_5$ free.
Converting back to equations and isolating basic variables:
$R_1: x_1 + x_5 = 18 \implies x_1 = 18 - x_5$
$R_2: x_2 - 2x_5 = 7 \implies x_2 = 7 + 2x_5$
$R_3: x_4 + 7x_5 = 2 \implies x_4 = 2 - 7x_5$
\\
\textbf{Answer:} C. $x_1 = 18 - x_5$, $x_2 = 7 + 2x_5$, $x_3$ is free, $x_4 = 2 - 7x_5$, $x_5$ is free.


\subsection*{Problem 1.2.14}
\textbf{Question:} Find the general solution of the system whose augmented matrix is given below.
\[ \begin{bmatrix} 1 & 0 & -5 & 0 & -8 & 3 \\ 0 & 1 & 2 & -1 & 0 & 4 \\ 0 & 0 & 0 & 0 & 1 & 0 \\ 0 & 0 & 0 & 0 & 0 & 0 \end{bmatrix} \]

\textbf{Solution:}
The pivots are in columns 1, 2, and 5. We must clear the non-zero entry above the pivot in column 5 (the $-8$ in $R_1$).
\[ \underrightarrow{R_1 \leftarrow R_1 + 8R_3} \begin{bmatrix} 1 & 0 & -5 & 0 & 0 & 3 \\ 0 & 1 & 2 & -1 & 0 & 4 \\ 0 & 0 & 0 & 0 & 1 & 0 \\ 0 & 0 & 0 & 0 & 0 & 0 \end{bmatrix} \]
The matrix is now in reduced echelon form.
Basic variables: $x_1, x_2, x_5$. Free variables: $x_3, x_4$.
Converting to equations and isolating basic variables:
$R_1: x_1 - 5x_3 = 3 \implies x_1 = 3 + 5x_3$
$R_2: x_2 + 2x_3 - x_4 = 4 \implies x_2 = 4 - 2x_3 + x_4$
$R_3: x_5 = 0$
\\
\textbf{Answer:} B. $x_1 = 3 + 5x_3$, $x_2 = 4 - 2x_3 + x_4$, $x_3$ is free, $x_4$ is free, $x_5 = 0$.


\subsection*{Problem 1.2.22}
\textbf{Question:} Determine the value(s) of $h$ such that the matrix is the augmented matrix of a consistent linear system.
\[ \begin{bmatrix} 1 & h & 3 \\ -5 & 20 & -18 \end{bmatrix} \]

\textbf{Solution:}
Row reduce the matrix to check for a row of the form $[0 \ 0 \ b]$.
\[ \underrightarrow{R_2 \leftarrow R_2 + 5R_1} \begin{bmatrix} 1 & h & 3 \\ 0 & 20 + 5h & -3 \end{bmatrix} \]
For the system to be inconsistent, the second row must equate to $0x_1 + 0x_2 = -3$. This happens if $20 + 5h = 0$.
Solving for $h$: $5h = -20 \implies h = -4$.
Therefore, if $h = -4$, the system is inconsistent. For the system to be \textit{consistent}, $h$ must be any number other than $-4$.
\\
\textbf{Answer:} A. The matrix is the augmented matrix of a consistent linear system if $h \neq -4$.


\subsection*{Problem 1.2.23}
\textbf{Question:} Choose $h$ and $k$ such that the system has (a) no solution, (b) a unique solution, and (c) many solutions.
\begin{align*}
    x_1 + hx_2 &= 4 \\
    5x_1 + 10x_2 &= k
\end{align*}

\textbf{Solution:}
Set up the augmented matrix and reduce it:
\[ \begin{bmatrix} 1 & h & 4 \\ 5 & 10 & k \end{bmatrix} \]
\[ \underrightarrow{R_2 \leftarrow R_2 - 5R_1} \begin{bmatrix} 1 & h & 4 \\ 0 & 10 - 5h & k - 20 \end{bmatrix} \]
\begin{itemize}
    \item \textbf{(a) No solution:} We need the last row to be a contradiction. This requires the coefficient of $x_2$ to be $0$ and the augmented part to be non-zero. 
    $10 - 5h = 0 \implies h = 2$.
    $k - 20 \neq 0 \implies k \neq 20$.
    \item \textbf{(b) Unique solution:} We need no free variables, which means every column in the coefficient matrix must have a pivot. We need the coefficient of $x_2$ to NOT be zero.
    $10 - 5h \neq 0 \implies h \neq 2$. As long as $h \neq 2$, $k$ can be any real number because consistency is guaranteed.
    \item \textbf{(c) Many solutions:} We need a free variable (meaning $10 - 5h = 0 \implies h = 2$) AND the system must still be consistent (meaning we cannot have a contradiction, so $k - 20 = 0 \implies k = 20$).
\end{itemize}
\textbf{Answer for Part 2 (b):} H. The system has a unique solution only when $h \neq 2$ and $k$ is any real number.


\subsection*{Problem 1.2.25}
\textbf{Question:} Determine whether the statement below is true or false. Justify the answer.
\textit{In some cases, a matrix may be row reduced to more than one matrix in reduced echelon form, using different sequences of row operations.}

\textbf{Solution:}
This tests the fundamental Uniqueness Theorem of Reduced Echelon Forms. While a matrix can be reduced to many different (unreduced) echelon forms depending on the choices you make during the forward phase, there is absolutely only \textit{one} correct reduced echelon form for any given matrix.
\\
\textbf{Answer:} D. The statement is false. Each matrix is row equivalent to one and only one reduced echelon matrix.


\subsection*{Problem 1.2.26}
\textbf{Question:} Determine whether the statement below is true or false. Justify the answer.
\textit{The echelon form of a matrix is unique.}

\textbf{Solution:}
This is the inverse of the previous question. Echelon forms are NOT unique. For instance, scaling a non-zero row in an echelon matrix produces a different matrix that is still in echelon form.
\\
\textbf{Answer:} B. The statement is false. The echelon form of a matrix is not unique, but the reduced echelon form is unique.


\subsection*{Problem 1.2.35}
\textbf{Question:} Suppose a $3 \times 6$ coefficient matrix for a system has three pivot columns. Is the system consistent? Why or why not?

\textbf{Solution:}
The coefficient matrix has $3$ rows. If it has $3$ pivot columns, that means there is a pivot in every single row. 
When we form the augmented matrix, we add a $7\text{th}$ column to represent the constants. 
Because every row already has a pivot in the coefficient portion, it is impossible for the augmented column (column 7) to be a pivot column (a row can only have one leading entry). 
Since the rightmost column cannot be a pivot column, the system cannot have a contradiction row $[0 \ 0 \ \dots \ 0 \ b]$. Therefore, it is always consistent.
\\
\textbf{Answer:} B. There is a pivot position in each row of the coefficient matrix. The augmented matrix will have seven columns and will not have a row of the form $[0 \ 0 \ 0 \ 0 \ 0 \ 0 \ 1]$, so the system is consistent.


\subsection*{Problem 1.2.36}
\textbf{Question:} Suppose a system of linear equations has a $3 \times 5$ augmented matrix whose fifth column is not a pivot column. Is the system consistent? Why or why not?

\textbf{Solution:}
\begin{itemize}
    \item \textbf{Part 1:} A linear system is consistent if and only if the rightmost column of the augmented matrix \textbf{is not} a pivot column. That is, if and only if an echelon form of the augmented matrix has \textbf{no row} of the form $[0 \ \dots \ 0 \ b]$ with $b$ nonzero.
    \item \textbf{Part 2:} In the augmented matrix described, is the rightmost column a pivot column? \textbf{No} (this is given directly in the problem description).
    \item \textbf{Part 4:} Therefore, by the Existence and Uniqueness Theorem, the linear system is \textbf{consistent}.
\end{itemize}


\subsection*{Problem 1.2.37}
\textbf{Question:} Suppose the coefficient matrix of a system of linear equations has a pivot position in every row. Explain why the system is consistent.

\textbf{Solution:}
This is a generalized version of problem 1.2.35. A row can only have one pivot position (it is the \textit{leading} entry). If every row in the coefficient matrix already possesses a leading entry, then when you append the augmented column, that augmented column can never contain a leading entry. Since the rightmost column is not a pivot column, the system is consistent.
\\
\textbf{Answer:} D. The system is consistent because the rightmost column of the augmented matrix is not a pivot column.

\newpage
\section{Part 3: In-Depth Analysis of Problems and Techniques}

\subsection{A) Problem Types and General Approach}

\textbf{1. Identifying Echelon Forms (Problem 1.2.2):}
\begin{itemize}
    \item \textit{Characteristics:} You are given matrices and asked to classify them.
    \item \textit{Approach:} Use a strict mental checklist. Are the zero rows at the bottom? (If no, $\to$ Neither). Do the leading non-zeros step to the right? (If no, $\to$ Neither). Are the leading entries $1$? (If no, $\to$ REF only). Are there zeros above the leading $1$s? (If no, $\to$ REF only. If yes to all, $\to$ RREF).
\end{itemize}

\textbf{2. Extracting General Solutions (Problems 1.2.11, 1.2.13, 1.2.14):}
\begin{itemize}
    \item \textit{Characteristics:} Given a matrix, you must write the solution parametrically.
    \item \textit{Approach:} 
    \begin{enumerate}
        \item Ensure the matrix is fully in RREF (backward phase complete).
        \item Identify Pivot Columns (Basic Variables).
        \item Identify Non-Pivot Columns (Free Variables).
        \item Write an equation for each non-zero row, isolating the basic variable on the left side, and pushing all free variables to the right side.
        \item Explicitly declare the free variables as "free".
    \end{enumerate}
\end{itemize}

\textbf{3. Parameter Analysis for Consistency (Problems 1.2.22, 1.2.23):}
\begin{itemize}
    \item \textit{Characteristics:} Matrix contains unknown constants $h, k$. 
    \item \textit{Approach:} Row reduce keeping $h$ and $k$ as algebraic expressions. To find conditions for inconsistency, set the coefficient expression to $0$ and the augmented expression to $\neq 0$. To find conditions for unique solutions, ensure every variable column has a non-zero coefficient (no free variables).
\end{itemize}

\textbf{4. Theoretical Application (Problems 1.2.25, 1.2.26, 1.2.35, 1.2.36, 1.2.37):}
\begin{itemize}
    \item \textit{Characteristics:} Questions asking about the implications of pivots.
    \item \textit{Approach:} These questions all test one core realization: \textbf{Consistency is determined entirely by the location of the pivots.} If a pivot lands in the last column, the system crashes. If the coefficient matrix claims all available rows for its pivots, the last column is safe, and the system is consistent.
\end{itemize}

\subsection{B) Key Algebraic and Calculus Manipulations}
\begin{itemize}
    \item \textbf{The Backward Phase (RREF):} In Problem 1.2.13, the matrix was given in echelon form but not reduced. The key manipulation was moving strictly from the \textit{right-most} pivot and working \textit{upwards} to clear the column, then moving to the next pivot to the left. If you try to clear the columns left-to-right during the backward phase, you will undo the zeros you just created.
    \item \textbf{Translating Zeros:} In Problem 1.2.11, reducing the matrix resulted in two rows of complete zeros. The manipulation trick is simply recognizing that $0 = 0$ provides no constraint, meaning those variables are free to be anything, leading directly to the parametric solution form.
\end{itemize}

\newpage
\section{Part 4: "Cheatsheet" and Tips for Success}

\subsection{Summary of the Algorithm}
\begin{enumerate}
    \item \textbf{Forward Phase:} (Yields REF)
    \begin{itemize}
        \item Find the leftmost nonzero column.
        \item Select a nonzero entry as a pivot. Swap rows if necessary to bring it to the top.
        \item Use row addition/subtraction to create zeros \textit{below} the pivot.
        \item Cover the row containing the pivot and ignore it. Repeat the process on the submatrix below.
    \end{itemize}
    \item \textbf{Backward Phase:} (Yields RREF)
    \begin{itemize}
        \item Begin with the rightmost pivot. Scale the row to make the pivot $1$.
        \item Use row addition/subtraction to create zeros \textit{above} the pivot.
        \item Move to the next pivot to the left and repeat.
    \end{itemize}
\end{enumerate}

\subsection{Tips and Shortcuts}
\begin{itemize}
    \item \textbf{How to easily spot a free variable:} Once the matrix is in REF, just look at the columns of the coefficient matrix. If a column does not have a "staircase corner" (leading entry) in it, the variable corresponding to that column is free.
    \item \textbf{Writing Solutions:} When writing parametric solutions, the basic variables are \textit{always} expressed in terms of the free variables. Never write a free variable in terms of a basic variable.
\end{itemize}

\subsection{Common Pitfalls}
\begin{itemize}
    \item \textbf{Stopping at REF:} When asked for a general solution, students often stop at row echelon form and try to do substitution algebraically. While technically possible, it is highly prone to error. Take the extra 30 seconds to run the backward phase to get RREF; the equations write themselves.
    \item \textbf{Misinterpreting a zero column:} If column 2 is all zeros, it means $x_2$ is a free variable (unless the system is inconsistent). It does \textit{not} mean $x_2 = 0$.
\end{itemize}

\newpage
\section{Part 5: Conceptual Synthesis and The "Big Picture"}

\subsection{A) Thematic Connections}
This section transitions from the arithmetic of Chapter 1.1 into the fundamental theory of Linear Algebra. The central theme here is \textbf{degrees of freedom}. In physics and engineering, the number of free variables directly corresponds to the number of independent dimensions a system can move in. By finding the reduced echelon form, we are essentially filtering out redundant information (basic variables) to find the core, unconstrained parameters (free variables) that define the system's behavior.

\subsection{B) Forward and Backward Links}
\begin{itemize}
    \item \textbf{Backward Link:} The row operations here are identical to 1.1, but now we apply them with a rigid, algorithmic discipline to guarantee a unique final form.
    \item \textbf{Forward Link:} The concepts of pivots and free variables are arguably the most important in the entire course. In Chapter 1.3-1.5, free variables will define the "span" of vectors. In Chapter 2, a matrix has an inverse if and only if it has no free variables. In Chapter 4, the number of pivot columns equals the "dimension" of the column space, and the number of free variables equals the dimension of the "null space."
\end{itemize}

\newpage
\section{Part 6: Real-World Application and Modeling}

\subsection{A) Concrete Scenarios}
\begin{enumerate}
    \item \textbf{Traffic Flow in Urban Networks:} Civil engineers use linear systems to model traffic flow through city intersections. The equations represent the law of conservation (vehicles entering an intersection must equal vehicles leaving). Because the network usually has more streets than intersections, the system has infinitely many solutions (free variables), which represent the maximum and minimum flow bounds for specific streets.
    \item \textbf{Chemical Equation Balancing:} In complex chemical reactions (e.g., combustion of a large hydrocarbon), balancing the equation by inspection is nearly impossible. Chemists set up a matrix where columns are molecules and rows are elements. Solving the matrix yields free variables, which are then scaled to the lowest integers to find the balanced stoichiometric coefficients.
\end{enumerate}

\subsection{B) Model Problem Setup: Traffic Flow}
Imagine a simple network of four one-way streets forming a square loop, with four intersections (A, B, C, D). 
Cars enter intersection A at a rate of 100 per hour and leave at intersection C at 100 per hour.
Let $x_1$ be the traffic flow from A to B.
Let $x_2$ be the flow from B to C.
Let $x_3$ be the flow from C to D.
Let $x_4$ be the flow from D to A.

At each intersection, Flow In = Flow Out.
Node A: $100 + x_4 = x_1 \implies x_1 - x_4 = 100$
Node B: $x_1 = x_2 \implies x_1 - x_2 = 0$
Node C: $x_2 = 100 + x_3 \implies x_2 - x_3 = 100$
Node D: $x_3 = x_4 \implies x_3 - x_4 = 0$

The augmented matrix to solve this model is:
\[ \begin{bmatrix} 1 & 0 & 0 & -1 & 100 \\ 1 & -1 & 0 & 0 & 0 \\ 0 & 1 & -1 & 0 & 100 \\ 0 & 0 & 1 & -1 & 0 \end{bmatrix} \]
Row reducing this matrix will yield a free variable, allowing the engineer to determine the maximum traffic capacity of street $x_1$ given constraints on the other streets.

\newpage
\section{Part 7: Common Variations and Untested Concepts}

A common problem type missing from this assignment is \textbf{Polynomial Interpolation}. This demonstrates how linear systems guarantee unique solutions for geometric curve fitting.

\subsection{Variation: Polynomial Interpolation}
\textbf{Concept:} Find a polynomial of degree 2 ($p(t) = a_0 + a_1t + a_2t^2$) whose graph passes through the specific points $(1, 12)$, $(2, 15)$, and $(3, 16)$.

\textbf{Setup and Solution:}
We plug the $(t, p(t))$ coordinates into the polynomial equation to create a system of linear equations where our unknowns are the coefficients $a_0, a_1, a_2$.
For $(1, 12): a_0 + a_1(1) + a_2(1)^2 = 12 \implies a_0 + a_1 + a_2 = 12$
For $(2, 15): a_0 + a_1(2) + a_2(2)^2 = 15 \implies a_0 + 2a_1 + 4a_2 = 15$
For $(3, 16): a_0 + a_1(3) + a_2(3)^2 = 16 \implies a_0 + 3a_1 + 9a_2 = 16$

The augmented matrix is:
\[ \begin{bmatrix} 1 & 1 & 1 & 12 \\ 1 & 2 & 4 & 15 \\ 1 & 3 & 9 & 16 \end{bmatrix} \]
Because the "time" points ($t=1, 2, 3$) are distinct, a theorem guarantees this matrix will row reduce to a unique solution (no free variables, no inconsistencies).
Reducing this matrix (via $R_2 - R_1, R_3 - R_1$, etc.) yields the unique solution:
$a_0 = 7, a_1 = 6, a_2 = -1$.
Thus, the unique polynomial is $p(t) = 7 + 6t - t^2$.

\newpage
\section{Part 8: Advanced Diagnostic Testing: "Find the Flaw"}

Below are 5 problems with complete, seemingly correct solutions. However, each contains a single fatal flaw. Identify the error and provide the correct calculation.

\hrule
\vspace{0.5cm}

\textbf{Flawed Problem 1:} Determine if the matrix is in reduced echelon form.
\[ \begin{bmatrix} 1 & 0 & -2 & 0 \\ 0 & 1 & 4 & 0 \\ 0 & 0 & 0 & 1 \end{bmatrix} \]
\textit{Flawed Solution:} The leading entries of the non-zero rows are $1$. They step down to the right. There are zeros below all the leading $1$s, and zeros above all the leading $1$s. All zero rows (there are none) are at the bottom. Therefore, it is in reduced echelon form.

\vspace{0.3cm}
\textbf{Flawed Problem 2:} Find the general solution to the system:
\[ \begin{bmatrix} 1 & 0 & 3 & 5 \\ 0 & 1 & -2 & 4 \\ 0 & 0 & 0 & 0 \end{bmatrix} \]
\textit{Flawed Solution:} 
The matrix is in RREF. The pivots are in columns 1 and 2. Column 3 has no pivot.
$x_1 + 3x_3 = 5 \implies x_3 = (5 - x_1)/3$
$x_2 - 2x_3 = 4 \implies x_2 = 4 + 2x_3$
Final Answer: $x_1$ is free, $x_2 = 4 + 2x_3$, $x_3 = (5 - x_1)/3$.

\vspace{0.3cm}
\textbf{Flawed Problem 3:} Determine the values of $h$ for which the system is consistent.
\[ \begin{bmatrix} 1 & -2 & h \\ 3 & -6 & 5 \end{bmatrix} \]
\textit{Flawed Solution:}
Replace $R_2$ with $R_2 - 3R_1$.
$R_2 \to \begin{bmatrix} 0 & 0 & 5 - 3h \end{bmatrix}$.
The matrix is $\begin{bmatrix} 1 & -2 & h \\ 0 & 0 & 5 - 3h \end{bmatrix}$.
Because there is a row of zeros on the left, the system has free variables and is therefore always consistent, regardless of the value of $h$.

\vspace{0.3cm}
\textbf{Flawed Problem 4:} Row reduce to RREF.
\[ \begin{bmatrix} 2 & 4 & 6 \\ 0 & 1 & 5 \end{bmatrix} \]
\textit{Flawed Solution:}
First, clear the entry above the pivot in column 2.
$R_1 \leftarrow R_1 - 4R_2$:
$R_1 \to \begin{bmatrix} 2 - 0 & 4 - 4 & 6 - 20 \end{bmatrix} = \begin{bmatrix} 2 & 0 & -14 \end{bmatrix}$.
Matrix is $\begin{bmatrix} 2 & 0 & -14 \\ 0 & 1 & 5 \end{bmatrix}$. 
The matrix satisfies all echelon properties and zeros are above the pivots, so this is the RREF.

\vspace{0.3cm}
\textbf{Flawed Problem 5:} Write the general solution for the augmented matrix.
\[ \begin{bmatrix} 1 & 2 & 0 & 0 \\ 0 & 0 & 1 & 0 \\ 0 & 0 & 0 & 1 \end{bmatrix} \]
\textit{Flawed Solution:}
Pivots are in cols 1, 3, and 4. Col 2 is free.
$x_1 = -2x_2$
$x_2$ is free
$x_3 = 0$
$x_4 = 1$
Final Answer: $( -2x_2, x_2, 0, 1 )$.

\end{document}